\section{Билет 16. Метод барьерных функций}

\ticketq{Метод барьерных функций.}

\paragraph{Постановка.}
\[
  \varphi(x) \to \min , \qquad
  S = \{\, x \mid \psi_i(x) \le 0,\ i = \overline{1,m} \,\} ,
  \qquad
  T = \{\, x \mid \psi_i(x) < 0 \,\} \ne \varnothing
\]
($T$ — внутренность $S$). Метод внутренних штрафов: итерации всё время
остаются строго внутри $S$, а к цели добавляется барьер, растущий до
бесконечности у границы.

\textbf{Важно:} метод неприменим к ограничениям-\emph{равенствам} — у
множества $\chi(x) = 0$ внутренности нет вовсе. Равенства обрабатывают внешним
штрафом (билет~15), барьером — только неравенства.

\paragraph{Барьерная функция.}
$B(x) \to +\infty$ при $x \to \partial S$ изнутри. Два стандартных вида
($\xi_i > 0$):
\[
  B(x) = -\sum_{i=1}^{m}\xi_i\,\frac{1}{\psi_i(x)} ,
  \qquad
  B(x) = -\sum_{i=1}^{m}\xi_i\,\ln\big(-\psi_i(x)\big) .
\]
Обратный барьер неудобен для дифференцирования, логарифмический удобнее.

\textbf{Почему барьер работает.} $\psi_i(x) \to 0-$ $\Rightarrow$
$-\psi_i(x) \to 0+$ $\Rightarrow$ $\ln(-\psi_i) \to -\infty$ $\Rightarrow$
$-\ln(-\psi_i) \to +\infty$ (аналогично $-1/\psi_i \to +\infty$). Кроме того,
$B$ определена \emph{только} на $T$, поэтому всякая точка с конечным
значением $\varphi + \mu B$ автоматически строго допустима — проверять
допустимость отдельно не нужно.

\paragraph{Вспомогательная задача.}
\[
  \varphi(x) + \mu B(x) \to \min_{x \in T} , \quad \mu > 0 ,
  \qquad x_\mu \ \text{— решение} ,
  \qquad \Theta(\mu) = \varphi(x_\mu) + \mu B(x_\mu) ,
\]
причём $\mu \to +0$ (в отличие от внешнего штрафа, где $\mu \to +\infty$).
Точки $x_\mu$ образуют \textbf{центральную траекторию}.

\paragraph{Лемма (свойства вспомогательных задач).}
\textit{Считаем $B \ge 0$ на $T$. Тогда}
\begin{csteps}
  \item $\min\limits_{S}\varphi \ \le\ \Theta(\mu)$ при всех $\mu > 0$;
  \item $B(x_\mu)$ не убывает при $\mu \to 0$
        ($\mu_1 > \mu_2 \Rightarrow B(x_{\mu_1}) \le B(x_{\mu_2})$);
  \item $\varphi(x_\mu)$ не возрастает при $\mu \to 0$;
  \item $\Theta(\mu)$ — неубывающая функция $\mu$;
  \item $\inf\limits_{\mu>0}\Theta(\mu) = \lim\limits_{\mu\to0}\Theta(\mu)$.
\end{csteps}

\textsc{Док-во.} Пусть $\mu_1 > \mu_2 > 0$, $x_i = x_{\mu_i}$.

\textbf{1)} $x_\mu \in T \subset S$ и $B \ge 0$ $\Rightarrow$
$\Theta(\mu) \ge \varphi(x_\mu) \ge \min_S\varphi$.

\textbf{2)} Из минимальности $x_1$ и $x_2$:
\[
  \varphi(x_1) + \mu_1 B(x_1) \le \varphi(x_2) + \mu_1 B(x_2) , \qquad
  \varphi(x_2) + \mu_2 B(x_2) \le \varphi(x_1) + \mu_2 B(x_1) .
\]
Складывая: $(\mu_1 - \mu_2)\big(B(x_1) - B(x_2)\big) \le 0$ $\Rightarrow$
$B(x_1) \le B(x_2)$.

\textbf{3)} Из второго неравенства:
$\varphi(x_2) - \varphi(x_1) \le \mu_2\big(B(x_1)-B(x_2)\big) \le 0$.

\textbf{4)}
$\Theta(\mu_2) \le \varphi(x_1) + \mu_2 B(x_1)
 \le \varphi(x_1) + \mu_1 B(x_1) = \Theta(\mu_1)$.

\textbf{5)} $\Theta$ монотонна и ограничена снизу (п.~1) $\Rightarrow$
нижняя грань достигается в пределе. $\blacksquare$

\paragraph{Теорема о сходимости.}
\textit{Пусть $T \ne \varnothing$, функции непрерывны на $T$, при каждом
$\mu > 0$ решение $x_\mu$ существует, $\{x_\mu\}$ лежит в компакте и всякая
точка $S$ приближается точками $T$ (то есть $\forall x \in S$,
$\forall\eps > 0$ $\exists y \in T$: $\norm{y - x} < \eps$). Тогда предельные
точки $\{x_\mu\}$ — решения задачи, причём}
\[
  \mu B(x_\mu) \to 0 , \qquad
  \min_S \varphi = \lim_{\mu\to0}\Theta(\mu) .
\]

\textsc{Док-во.} По условию для любого $\eps$ найдётся $y \in T$ с
$\varphi(y) < \varphi_* + \eps$, где $\varphi_* = \min_S\varphi$. Тогда
\[
  \Theta(\mu) \le \varphi(y) + \mu B(y) \xrightarrow[\mu\to0]{} \varphi(y)
  \ \Rightarrow\ \lim_{\mu\to0}\Theta(\mu) \le \varphi_* + \eps
  \quad \forall\eps
  \ \Rightarrow\ \lim\Theta \le \varphi_* ,
\]
а по п.~1 леммы $\lim\Theta \ge \varphi_*$. Значит
$\lim_{\mu\to0}\Theta(\mu) = \varphi_*$. Далее
\[
  0 \le \mu B(x_\mu) = \Theta(\mu) - \varphi(x_\mu)
    \le \Theta(\mu) - \varphi_* \to 0 .
\]
Для предельной точки $\bar{x}$: $\bar{x} \in S$ (замкнутость) и
$\varphi(\bar{x}) = \lim\varphi(x_\mu) \le \lim\Theta = \varphi_*$, то есть
$\bar{x}$ — решение. $\blacksquare$

\paragraph{Алгоритм.}
\begin{csteps}
  \item Задать $x_1 \in T$ (строго допустимую), $\mu_1 > 0$, $\eps > 0$,
        $\beta \in (0,1)$, $k = 1$.
  \item Найти $x_{\mu_k} = \argmin_{x\in T}\big[\varphi(x) + \mu_k B(x)\big]$,
        стартуя с предыдущей точки.
  \item Если $\mu_k B(x_{\mu_k}) < \eps$ — конец.
  \item Иначе $\mu_{k+1} = \beta\mu_k$, $k := k+1$, к~шагу~2.
\end{csteps}

\paragraph{Ограничения.}
Нужна строго допустимая $x_1$; метод неприменим при $T = \varnothing$; точки
не могут лежать на границе, где барьер не определён, — а оптимум, как правило,
именно там; при малых $\mu$ функция $\varphi + \mu B$ плохо обусловлена у
границы; равенства напрямую не поддерживаются.

\seefull{гл.~3, \S3 и \S12 (численный пример в \S12)}
